\part{The Unquiet Dead}

\section{The Spirit Taxonomy}

The dead do not remain for a single reason, nor do they manifest in a single form. Over years of investigation---and, I confess, more than a few exorcisms conducted in haste---I have developed a working taxonomy of the unquiet. This is not an academic exercise. The method of laying a spirit to rest depends entirely upon \emph{why} it remains and \emph{how} it manifests.

\subsection{Class I: The Echo}

The Echo is not truly a spirit. It is a \emph{recording}---a psychic imprint left by a traumatic death, repeated endlessly until the conditions that created it are addressed. Echoes do not speak, do not learn, and cannot be negotiated with. They are the scream that plays again at midnight, the footsteps that cross the same corridor at the same hour, the face that appears in the window of a house that burned a century ago.

\textbf{Investigation Notes:} Echoes require witness to fade. A single observer who speaks aloud the truth of the event---who names the dead and describes their end---will often break the loop. If the death was unjust, the Echo may persist until justice is done. In such cases, I have found that a public acknowledgment of the wrong, witnessed by a representative of the community, suffices where private ritual fails.

\textbf{The Witness's Technique:} When I encounter an Echo, I place a small silvered mirror facing the site of the imprint. The Echo will play once into the glass. Afterward, the mirror may be shattered---each shard holds a fragment of the memory, which can be used in later investigations or ritually cleansed.

\textbf{Case Study: \emph{The Lamplighter of the Third Stack}} (Aelerian Ranges, 847). A miner's death in a gas explosion repeated every night at the same hour. The Echo was not the man's death, but his final attempt to ring the warning bell. We placed a recording mirror at the gallery entrance, then had the mine's foreman read the names of all who had died in that disaster into the glass. The Echo ceased. The foreman slept peacefully for the first time in thirty years.

\subsection{Class II: The Anchor}

An Anchor is a spirit bound to a specific object, location, or person by unresolved emotion---grief, rage, love, or duty. Unlike the Echo, an Anchor possesses will and memory. It can speak, bargain, and choose to manifest or withhold itself. It is, in essence, a person who has forgotten how to leave.

\textbf{Variants by Cultural Tradition:}

\begin{itemize}
  \item \emph{The Hearth-Bound (Vilikari, Aelaerem):} Spirits of grandparents or parents who refuse to depart until they have ensured the flourishing of their descendants. These Anchors are generally benevolent, though they may become irritable if their advice is ignored repeatedly.
  \item \emph{The Duty-Bound (Aeler, Kahfagia):} Soldiers, watchmen, and sailors who died at their post and continue to perform their duty. I once met a lighthouse keeper who had been dead for seven years. The lamp never went out. He did not speak, but his logbook was written in a hand that was not quite attached to a wrist.
  \item \emph{The Wronged (Universally encountered):} Victims of murder, betrayal, or accident who refuse to pass until their death is acknowledged and, where possible, avenged. These are the most dangerous Anchors, as their grief curdles into something that resembles hunger.
\end{itemize}

\textbf{Investigation Notes:} An Anchor cannot be exorcised by force without risk of creating a Poltergeist (Class IV). The proper method is negotiation. Determine what the spirit requires to complete its business. Offer witness---a third party who will attest to the resolution. In cases where the spirit's demand is impossible (the murderer died decades ago, the lover has since remarried), a ritual of substitution may suffice.

\textbf{The Witness's Technique:} I conduct Anchor interrogations in a room with three mirrors---one facing the spirit, one facing myself, and one placed between us to catch the reflection of the truth. The Anchor cannot lie in the space between two reflections. It may refuse to speak, but it cannot deceive.

\textbf{Case Study: \emph{The Widow of the Red Tower}} (Vhasia, 852). A noblewoman's ghost haunted her husband's estate, refusing to pass until he confessed to poisoning her. I placed three mirrors in the great hall and asked the spirit to show me her death. The reflection did not lie. The husband confessed before dawn. The widow vanished with a smile that I still see in my dreams.

\subsection{Class III: The Procession}

A Procession is not a single spirit, but a \emph{route}---a path that the dead walk, sometimes for centuries, between two points of significance. Processions are most common in regions where death was sudden and widespread: battlefields, plague pits, shipwreck coasts.

\textbf{Characteristics:}
\begin{itemize}
  \item Processions occur on a regular schedule (nightly, monthly, on the anniversary of the disaster).
  \item They follow a fixed route that predates the current roads.
  \item The spirits within a Procession do not interact with observers. They are \emph{marching}, not haunting.
  \item Interrupting a Procession---by standing in its path, attempting to speak to its members, or building over its route---provokes a violent response from the entire column.
\end{itemize}

\textbf{Investigation Notes:} Processions are not malevolent, but they are \emph{inflexible}. The proper response is to observe, document, and---if the route must be altered---negotiate with the \emph{destination}, not the dead. If a Procession marches to a specific grave or shrine, moving that terminus will often reroute the spirits.

\textbf{The Traveler's Intervention:} When I need to cross where a Procession walks, I invoke The Traveler. A single pinch of road dust from a crossroads where three murders occurred, blown toward the column, causes the spirits to \emph{step aside} for exactly the time needed to pass. I have used this technique four times. On three occasions, it worked. On the fourth, the Procession followed me for a week. I do not recommend relying on this method unless you are prepared to run.

\textbf{Case Study: \emph{The Phantom Caravan of the Saffron Pass}} (Dhahara, 855). A merchant's new road cut directly across the path of a weekly Procession. The column attacked every merchant who attempted the crossing. I consulted with a Fhara caravaneer, who revealed that the Procession was marching from an old well to a date grove. We moved a single stone---a marker that had served as the terminus---to the side of the new road. The Procession followed the stone. The road was safe by nightfall.

\subsection{Class IV: The Poltergeist}

The Poltergeist is an Anchor that has been corrupted by rage, starvation, or the failure of negotiation. It is no longer a person---it is a \emph{wound} that has learned to strike back. Poltergeists manifest as telekinetic force, localized environmental disturbances, and, in advanced cases, violent possession of the living.

\textbf{Stages of Poltergeist Development:}
\begin{description}
  \item[Stage One (The Knock):] Unexplained sounds, objects shifting slightly between observations, drafts in sealed rooms.
  \item[Stage Two (The Grasp):] Objects thrown with force, doors slamming, furniture overturned. Living creatures experience nightmares, nausea, and a persistent sense of being watched.
  \item[Stage Three (The Hunger):] The Poltergeist begins to feed on fear and pain. It will target specific individuals, escalating from psychological torment to physical harm. Bruises, scratches, and broken bones become common.
  \item[Stage Four (The Mask):] The Poltergeist possesses a living host, using their body to commit acts of violence. Exorcism at this stage is dangerous and often fatal to the host.
\end{description}

\textbf{Exorcism with The Sentinel:} The Sentinel's power is that of the witnessed boundary. I perform the following ritual:
\begin{enumerate}
  \item Identify the Anchor's original location of death or burial. This is the \emph{root}.
  \item Surround the root with a circle of salt and iron filings, blessed by a representative of the local faith. (The Sentinel does not care which faith; only that the witness is genuine.)
  \item Place a silvered mirror at each cardinal point within the circle, facing outward.
  \item Speak the name of the dead three times---once for acknowledgment, once for judgment, once for release.
  \item Shatter the mirrors simultaneously. The Sentinel seals the break.
\end{enumerate}
This ritual does not \emph{destroy} the Poltergeist. It returns it to the state of an Anchor, at which point standard negotiation may resume.

\textbf{Case Study: \emph{The Plague House of Port Shed}} (Dhaharan Coast, 858). A family of Sidhi dockworkers was being terrorized by a Poltergeist in Stage Three. The Anchor was a young woman who had died of a fever three months prior---not violently, but \emph{alone}, her family having quarantined her. I performed The Sentinel's containment. At the moment of release, I placed a small hand mirror in the center of the circle, reflecting her own face. The Poltergeist collapsed into a weeping Anchor. She passed into the care of a local medium. She has not returned.

